% Imporditud: tools/import_eksperimendid.py
% Allikas: Eesti füüsikaolümpiaadi eksperimendi ülesannete
%   kogu (Rael Kalda, 2023), ülesanne Ü28, lahendus L28.
\setAuthor{Erkki Tempel}
\setRound{piirkonnavoor}
\setYear{2020}
\setNumber{P E1}
\setDifficulty{1}
\setTopic{varia}
\setVahendid{spagetid, 1 makaron}
\setPunktid{10}
\setAllikas{Eksperimendi kogu Ü28, lahendus lk 44}
\setLahendusseis{toores}

\prob{Makaronid}
Hinnata, mitu korda erineb ühe spageti mass ühe makaroni massist.
\textit{Märkus:} Spagetti tohib tükkideks murda.

\hint

\solu
Kasutame spagetti kangina. (1 p) Täpsema tulemuse saamiseks tuleb kasutada kah-

te või kolme spagetti koos. (1 p)

Asetame torumakaroni spagettide otsa ning tasakaalustame laua serval. (1 p) Mõõ-

dame spagettide masskeskme ning torumakaroni masskeskme kauguse toetuspunk-

tist (jõuõlad). (1 p) Spagettide masskese asub spagettide keskel.

Mõõtmiseks saame kasutada väikest spageti juppi. (1 p) Leiame, mitu korda erine-

vad spageti ja makaroni jõuõlgade pikkused.

Kangireeglist msls = mmlm saame leida, mitu korda erinevad torumakaroni ja

spageti massid. mm = ls (3 p) ms lm

Kordusmõõtmised --- 1 p.

Õige tulemus täpsusega $\pm$20 \% --- 1 p.
\probend
